% ============================================================
\section{11.3 그래디언트 소실의 재발과 LSTM/GRU}
\begin{frame}{그래디언트 소실이 시간 축에서 재발}
  Week 9.2에서 본 것과 \textbf{정확히 같은 패턴}:
  \begin{itemize}
    \item $\tanh'$의 최댓값은 1이지만 \textbf{대부분의 구간에서 1보다
      작}고
    \item 여기에 $W_{hh}$까지 $T$번(시퀀스 길이만큼) 곱해진다
    \item 시퀀스가 길수록($T$가 클수록) 이 곱은
    \begin{itemize}
      \item $W_{hh}$ 고유값 $< 1$: 지수적으로 0에 가까워짐
        --- \kb{그래디언트 소실}
      \item $W_{hh}$ 고유값 $> 1$: \textbf{발산}
        --- \kb{그래디언트 폭발}
    \end{itemize}
  \end{itemize}
  \vskip0.5em
  \begin{itemize}
    \item 결과: 기본 RNN은 \textbf{먼 과거의 정보를 거의 기억하지
      못}한다 --- ``10단어 전에 나온 주어''를 지금 활용해야 하는
      문장에서 특히 취약
  \end{itemize}
  \vskip0.4em
  숫자로 ($\tanh' \approx 0.5$, $W_{hh}$ 크기 0.9, $T=20$):
  \[ (0.5 \times 0.9)^{20} = 0.45^{20} \approx 4.5 \times 10^{-7}
     \;\;(\text{사실상 0}) \]
  {\footnotesize 20단어 전의 정보는 지금 시점의 학습에 거의 아무
  영향도 주지 못한다.}
\end{frame}
